\section{Flat modules}
\label{section-flat}

\noindent
We can define flat modules exactly as in the case of modules over rings.

\begin{definition}
\label{definition-flat}
Let $(X, \mathcal{O}_X)$ be a ringed space.
An $\mathcal{O}_X$-module $\mathcal{F}$ is {\it flat} if the functor
$$
\textit{Mod}(\mathcal{O}_X)
\longrightarrow
\textit{Mod}(\mathcal{O}_X), \quad
\mathcal{G} \mapsto \mathcal{G} \otimes_{\mathcal{O}_X} \mathcal{F}
$$
is exact.
\end{definition}

\noindent
We can characterize flatness by looking at the stalks.

\begin{lemma}
\label{lemma-flat-stalks-flat}
Let $(X, \mathcal{O}_X)$ be a ringed space.
An $\mathcal{O}_X$-module $\mathcal{F}$ is flat
if and only if the stalk $\mathcal{F}_x$ is a flat
$\mathcal{O}_{X, x}$-module for all $x \in X$.
\end{lemma}

\begin{proof}
Assume $\mathcal{F}_x$ is a flat $\mathcal{O}_{X, x}$-module for all
$x \in X$. In this case, if $\mathcal{G} \to \mathcal{H} \to \mathcal{K}$
is exact, then also
$\mathcal{G} \otimes_{\mathcal{O}_X} \mathcal{F} \to
\mathcal{H} \otimes_{\mathcal{O}_X} \mathcal{F} \to
\mathcal{K} \otimes_{\mathcal{O}_X} \mathcal{F}$
is exact because we can check exactness at stalks and because
tensor product commutes with taking stalks, see
Lemma \ref{lemma-stalk-tensor-product}.
Conversely, suppose that $\mathcal{F}$ is flat, and let $x \in X$.
Consider the skyscraper sheaves $i_{x, *} M$ where $M$ is a
$\mathcal{O}_{X, x}$-module. Note that
$$
M \otimes_{\mathcal{O}_{X, x}} \mathcal{F}_x =
\left(i_{x, *} M \otimes_{\mathcal{O}_X} \mathcal{F}\right)_x
$$
again by
Lemma \ref{lemma-stalk-tensor-product}.
Since $i_{x, *}$ is exact, we see that the fact that $\mathcal{F}$
is flat implies that $M \mapsto M \otimes_{\mathcal{O}_{X, x}} \mathcal{F}_x$
is exact. Hence $\mathcal{F}_x$ is a flat $\mathcal{O}_{X, x}$-module.
\end{proof}

\noindent
Thus the following definition makes sense.

\begin{definition}
\label{definition-flat-at-point}
Let $(X, \mathcal{O}_X)$ be a ringed space. Let $x \in X$.
An $\mathcal{O}_X$-module $\mathcal{F}$ is
{\it flat at $x$} if $\mathcal{F}_x$ is a flat
$\mathcal{O}_{X, x}$-module.
\end{definition}

\noindent
Hence we see that $\mathcal{F}$ is a flat $\mathcal{O}_X$-module
if and only if it is flat at every point.

\begin{lemma}
\label{lemma-base-change-flat}
Let $f : (X, \mathcal{O}_X) \to (Y, \mathcal{O}_Y)$
be a morphism of ringed spaces. If $\mathcal{G}$ is a flat
$\mathcal{O}_Y$-module, then $f^*\mathcal{G}$ is a flat
$\mathcal{O}_X$-module.
\end{lemma}

\begin{proof}
Combine Lemma \ref{lemma-flat-stalks-flat} with
Sheaves, Lemma \ref{sheaves-lemma-stalk-pullback-modules} and
Algebra, Lemma \ref{algebra-lemma-flat-base-change}.
\end{proof}

\begin{lemma}
\label{lemma-colimits-flat}
Let $(X, \mathcal{O}_X)$ be a ringed space.
A filtered colimit of flat $\mathcal{O}_X$-modules is flat.
A direct sum of flat $\mathcal{O}_X$-modules is flat.
\end{lemma}

\begin{proof}
This follows from
Lemma \ref{lemma-tensor-commute-colimits},
Lemma \ref{lemma-stalk-tensor-product},
Algebra, Lemma \ref{algebra-lemma-directed-colimit-exact},
and the fact that we can check exactness at stalks.
\end{proof}

\begin{lemma}
\label{lemma-j-shriek-flat}
Let $(X, \mathcal{O}_X)$ be a ringed space.
Let $U \subset X$ be open. The sheaf $j_{U!}\mathcal{O}_U$
is a flat sheaf of $\mathcal{O}_X$-modules.
\end{lemma}

\begin{proof}
The stalks of $j_{U!}\mathcal{O}_U$ are either zero or equal
to $\mathcal{O}_{X, x}$. Apply
Lemma \ref{lemma-flat-stalks-flat}.
\end{proof}

\begin{lemma}
\label{lemma-module-quotient-flat}
Let $(X, \mathcal{O}_X)$ be a ringed space.
\begin{enumerate}
\item Any sheaf of $\mathcal{O}_X$-modules is a quotient of
a direct sum $\bigoplus j_{U_i!}\mathcal{O}_{U_i}$.
\item Any $\mathcal{O}_X$-module is a quotient of
a flat $\mathcal{O}_X$-module.
\end{enumerate}
\end{lemma}

\begin{proof}
Let $\mathcal{F}$ be an $\mathcal{O}_X$-module.
For every open $U \subset X$ and every
$s \in \mathcal{F}(U)$ we get a morphism
$j_{U!}\mathcal{O}_U \to \mathcal{F}$, namely the adjoint to
the morphism $\mathcal{O}_U \to \mathcal{F}|_U$, $1 \mapsto s$.
Clearly the map
$$
\bigoplus\nolimits_{(U, s)} j_{U!}\mathcal{O}_U
\longrightarrow
\mathcal{F}
$$
is surjective, and the source is flat by combining Lemmas
\ref{lemma-colimits-flat} and \ref{lemma-j-shriek-flat}.
\end{proof}

\begin{lemma}
\label{lemma-flat-tor-zero}
Let $(X, \mathcal{O}_X)$ be a ringed space.
Let
$$
0 \to \mathcal{F}'' \to \mathcal{F}' \to \mathcal{F} \to 0
$$
be a short exact sequence of $\mathcal{O}_X$-modules.
Assume $\mathcal{F}$ is flat. Then for any $\mathcal{O}_X$-module
$\mathcal{G}$ the sequence
$$
0 \to
\mathcal{F}'' \otimes_\mathcal{O} \mathcal{G} \to
\mathcal{F}' \otimes_\mathcal{O} \mathcal{G} \to
\mathcal{F} \otimes_\mathcal{O} \mathcal{G} \to 0
$$
is exact.
\end{lemma}

\begin{proof}
Using that $\mathcal{F}_x$ is a flat $\mathcal{O}_{X, x}$-module
for every $x \in X$ and that exactness can be checked on stalks, this
follows from
Algebra, Lemma \ref{algebra-lemma-flat-tor-zero}.
\end{proof}

\begin{lemma}
\label{lemma-flat-ses}
\begin{slogan}
Kernels of epimorphisms and extensions of flat sheaves of modules over
a ringed space are again flat.
\end{slogan}
Let $(X, \mathcal{O}_X)$ be a ringed space.
Let
$$
0 \to
\mathcal{F}_2 \to
\mathcal{F}_1 \to
\mathcal{F}_0 \to 0
$$
be a short exact sequence of $\mathcal{O}_X$-modules.
\begin{enumerate}
\item If $\mathcal{F}_2$ and $\mathcal{F}_0$ are flat so is
$\mathcal{F}_1$.
\item If $\mathcal{F}_1$ and $\mathcal{F}_0$ are flat so is
$\mathcal{F}_2$.
\end{enumerate}
\end{lemma}

\begin{proof}
Since exactness and flatness may be checked at the level of stalks
this follows from
Algebra, Lemma \ref{algebra-lemma-flat-ses}.
\end{proof}

\begin{lemma}
\label{lemma-flat-resolution-of-flat}
Let $(X, \mathcal{O}_X)$ be a ringed space.
Let
$$
\ldots \to
\mathcal{F}_2 \to
\mathcal{F}_1 \to
\mathcal{F}_0 \to
\mathcal{Q} \to 0
$$
be an exact complex of $\mathcal{O}_X$-modules.
If $\mathcal{Q}$ and all $\mathcal{F}_i$ are flat $\mathcal{O}_X$-modules,
then for any $\mathcal{O}_X$-module $\mathcal{G}$ the complex
$$
\ldots \to
\mathcal{F}_2 \otimes_{\mathcal{O}_X} \mathcal{G} \to
\mathcal{F}_1 \otimes_{\mathcal{O}_X} \mathcal{G} \to
\mathcal{F}_0 \otimes_{\mathcal{O}_X} \mathcal{G} \to
\mathcal{Q} \otimes_{\mathcal{O}_X} \mathcal{G} \to 0
$$
is exact also.
\end{lemma}

\begin{proof}
Follows from Lemma \ref{lemma-flat-tor-zero} by splitting the complex
into short exact sequences and using Lemma \ref{lemma-flat-ses} to
prove inductively that $\Im(\mathcal{F}_{i + 1} \to \mathcal{F}_i)$
is flat.
\end{proof}

\noindent
The following lemma gives one direction of the equational criterion of
flatness (Algebra, Lemma \ref{algebra-lemma-flat-eq}).

\begin{lemma}
\label{lemma-flat-eq}
Let $(X, \mathcal{O}_X)$ be a ringed space. Let $\mathcal{F}$ be a flat
$\mathcal{O}_X$-module. Let $U \subset X$ be open and let
$$
\mathcal{O}_U \xrightarrow{(f_1, \ldots, f_n)}
\mathcal{O}_U^{\oplus n} \xrightarrow{(s_1, \ldots, s_n)}
\mathcal{F}|_U
$$
be a complex of $\mathcal{O}_U$-modules. For every $x \in U$ there
exists an open neighbourhood $V \subset U$ of $x$ and a factorization
$$
\mathcal{O}_V^{\oplus n}
\xrightarrow{A}
\mathcal{O}_V^{\oplus m} \xrightarrow{(t_1, \ldots, t_m)}
\mathcal{F}|_V
$$
of $(s_1, \ldots, s_n)|_V$ such that $A \circ (f_1, \ldots, f_n)|_V = 0$.
\end{lemma}

\begin{proof}
Let $\mathcal{I} \subset \mathcal{O}_U$ be the sheaf of ideals
generated by $f_1, \ldots, f_n$. Then $\sum f_i \otimes s_i$ is
a section of $\mathcal{I} \otimes_{\mathcal{O}_U} \mathcal{F}|_U$
which maps to zero in $\mathcal{F}|_U$. As $\mathcal{F}|_U$ is flat
the map
$\mathcal{I} \otimes_{\mathcal{O}_U} \mathcal{F}|_U \to \mathcal{F}|_U$
is injective. Since $\mathcal{I} \otimes_{\mathcal{O}_U} \mathcal{F}|_U$
is the sheaf associated to the presheaf tensor product, we see
there exists an open neighbourhood $V \subset U$ of $x$ such
that $\sum f_i|_V \otimes s_i|_V$ is zero in
$\mathcal{I}(V) \otimes_{\mathcal{O}(V)} \mathcal{F}(V)$.
Unwinding the definitions using Algebra, Lemma \ref{algebra-lemma-relations}
we find $t_1, \ldots, t_m \in \mathcal{F}(V)$ and $a_{ij} \in \mathcal{O}(V)$
such that $\sum a_{ij}f_i|_V = 0$ and $s_i|_V = \sum a_{ij}t_j$.
\end{proof}














